% !TeX program = xelatex
\documentclass[t,aspectratio=169,9pt]{beamer}
\makeatletter\def\input@path{{../../theme/}}\makeatother   % το θέμα είναι στο Presentations/theme/
\usetheme{pdp}

\title{Εξαντλητική εξέταση}
\subtitle{Επίπεδο Junior}
\author{Αναστάσης Κολιόπουλος}
% \institute{Πανελλήνιος Διαγωνισμός Πληροφορικής}
\date{}

\begin{document}
\pdptitleframe                 % πρώτη σελίδα
\pdptocframe{Περιεχόμενα}      % μπλε σελίδα με όλες τις ενότητες

\section{Εισαγωγή}        % πριν από κάθε ενότητα μπαίνει αυτόματα μπλε σελίδα περιεχομένων
% \begin{frame}{Τίτλος διαφάνειας}{Υπότιτλος (προαιρετικός)}
%   ...
% \end{frame}

\section{Επαναληπτική εξαντλητική εξέταση}

\subsection{Εγγύτερο ζεύγος σημείων}
\begin{frame}{Εγγύτερο ζεύγος σημείων}
    \field{Είσοδος:}
    \begin{itemize}
        \pause
        \item Σύνολο σημείων $P = \{(x_{1}, y_{1}), (x_{2}, y_{2}), \dots, (x_{n}, y_{n}) \}\subset \mathbb{R}^{2}$.
    \end{itemize}
    \pause
    \field{Έξοδος:} 
    \begin{itemize}
        \pause
        \item Η ελάχιστη Ευκλείδεια απόσταση μεταξύ δύο σημείων, δηλαδή, η εύρεση της τιμής 
        \begin{align*}
            \text{minimumDistance} = \min_{1 \leq i < j \leq n}\left\{\sqrt{(x_{i} - x_{j})^{2} + (y_{i} - y_{j})^{2}}\right\}.
        \end{align*}
    \end{itemize}
    \pause
    \field{Λύση:}
    \begin{itemize}
        \pause 
        \item Σύμφωνα με τεχνική εξάντλησης, «εξαντλούμε» όλα τα ζεύγη σημείων $(i, j)$. Πιο συγκεκριμένα, για κάθε ζεύγος
        υπολογίζουμε την απόσταση του $(x_i, y_i)$ από το $(x_j, y_j)$.
        \pause 
        \item Η ελάχιστη εξ' αυτών είναι η απάντηση στο πρόβλημα.
    \end{itemize}
\end{frame}

\begin{frame}[fragile]{Εγγύτερο ζεύγος σημείων}
\begin{cpp}
double minimumDistance = numeric_limits<double>::max();

for (int i = 0; i < N - 1; ++i) 
{
    for (int j = i + 1; j < N; ++j)
    {
        double currentDistance = sqrt((X[i] - X[j]) * (X[i] - X[j]) + (Y[i] - Y[j]) * (Y[i] - Y[j]));
        minimumDistance = min(minimumDistance, currentDistance);
    }
}
\end{cpp}
\pause
\vspace{0.25cm}
\field{Πολυπλοκότητα:}
\begin{itemize}
    \pause
    \item Δύο φωλιασμένοι βρόχοι, καθένας εκ των οποίων εκτελεί $\mathcal{O}(n)$ επαναλήψεις.
    \pause
    \item Υπολογισμός συνάρτησης \texttt{sqrt} εντός των βρόχων, πολυπλοκότητας $\mathcal{O}(\lg a)$. %, όπου 
    % \begin{align*}
    %     a = \max_{1 \leq i < j \leq n}\left\{(x_{i} - x_{j})^{2} + (y_{i} - y_{j})^{2}\right\}.
    % \end{align*}
    \pause
    \item Κάποιες πράξεις κόστους $\mathcal{O}(1)$ εντός των βρόχων.
\end{itemize}
\end{frame}

\begin{frame}[fragile]{Εγγύτερο ζεύγος σημείων}
    \begin{itemize}
        \item Συνολικά, ο χρόνος εκτέλεσης είναι $\mathcal{O}(n^2 \lg a)$, όπου 
        \begin{align*}
            a = \max_{1 \leq i < j \leq n}\{(x_{i} - x_{j})^{2} + (y_{i} - y_{j})^{2}\}.
        \end{align*}
    \end{itemize}
    \pause
    \begin{block}{Παρατήρηση.} 
        Δεδομένου ότι οι αποστάσεις είναι θετικές, η ρίζα δεν επηρεάζει τη μεταξύ τους
        σύκριση και άρα μπορούμε να την αγνοήσουμε, μέχρις ότου εκτυπώσουμε το αποτέλεσμα. Έτσι, 
        έχουμε ότι
        \begin{cpp}
            double currentDistance = (X[i] - X[j]) * (X[i] - X[j]) + (Y[i] - Y[j]) * (Y[i] - Y[j]);
        \end{cpp}
        και επίσης, στο τέλος προσθέτουμε την εντολή
        \begin{cpp}
            minimumDistance = sqrt(minimumDistance);
        \end{cpp}
    \end{block}
    %\vspace{0.25cm}
    \pause
    \field{Πολυπλοκότητα μετά την αλλαγή:}
    \begin{itemize}
        \pause
        \item $\mathcal{O}(n^2 + \lg a)$.
    \end{itemize}
\end{frame}

\subsection{Μέγιστο άθροισμα υποπίνακα}
\begin{frame}{Μέγιστο άθροισμα υποπίνακα}
    \field{Είσοδος:}
    \begin{itemize}
        \pause
        \item Πίνακας $A[1..n]$ ακέραιων αριθμών.
    \end{itemize}
    \pause
    \field{Έξοδος:} 
    \begin{itemize}
        \pause
        \item Το μέγιστο άθροισμα διαδοχικών αριθμών του πίνακα. Ισοδύναμα, 
        \begin{align*}
            \text{maximumSubarraySum} = \max_{1 \leq i \leq j \leq n}\left\{\sum_{k = i}^{j} A[k]\right\}.
        \end{align*}
    \end{itemize}
    \pause
    \field{Λύση:}
    \begin{itemize}
        \pause 
        \item Δοκιμάζουμε κάθε πιθανό υποπίνακα $A$. Πιο συγκεκριμένα, για κάθε δυνατή αρχή $i$, δοκιμάζουμε κάθε δυνατό τέλος $j$, υπολογίζοντας 
        το άθροισμα των στοιχείων $A[i], A[i + 1], \dots, A[j]$.
        \pause
        \item Το μεγαλύτερο από αυτά αποτελεί και την απάντηση.
    \end{itemize}
\end{frame}

\begin{frame}[fragile]{Μέγιστο άθροισμα υποπίνακα}
\begin{cpp}
int maximumSubarraySum = numeric_limits<int>::min();

for (int i = 1; i <= N; ++i)
{
    for (int j = i; j <= N; ++j)
    {
        int currentSum = 0;

        for (int k = i; k <= j; ++k)
        {
            currentSum += A[k];
        }

        maximumSubarraySum = max(maximumSubarraySum, currentSum);
    }
}
\end{cpp}
\end{frame}

\begin{frame}{Μέγιστο άθροισμα υποπίνακα}
Για $i = 3$, εξετάζουμε τους υποπίνακες ως εξής.

\vspace{0.5cm}
\centering
\begin{tikzpicture}
% Κελιά: η αρχή i σκούρο μπλε, τα δυνατά τέλη j ανοιχτό μπλε
\filldraw[pdpblue] (2,0) rectangle (3,1);
\filldraw[pdpblue!15] (3,0) rectangle (6,1);

\draw[pdpmuted] (0,0) grid (6,1);

% Όνομα του πίνακα
\node[anchor=east] at (-0.2,0.5) {$A$};

% Δείκτες
\foreach \k in {1,...,6}
  \node[text=pdpmuted, font=\small] at (\k - 0.5, 1.35) {$\k$};

% Τιμές
\node at (0.5,0.5) {$1$};
\node at (1.5,0.5) {$-3$};
\node[text=white] at (2.5,0.5) {$\mathbf{4}$};
\node at (3.5,0.5) {$-2$};
\node at (4.5,0.5) {$5$};
\node at (5.5,0.5) {$-1$};

% Η αρχή i: βέλος πάνω από το κελί 3, σταθερό σε όλα τα βήματα
\draw[pdpblue, thick, -latex] (2.5,2.35) -- (2.5,1.65);
\node[text=pdpblue, above] at (2.5,2.35) {$i$};

% Μία αγκύλη από την αρχή i = 3 που σε κάθε βήμα επεκτείνεται κατά ένα κελί,
% και ένα βέλος j που μετακινείται στο τρέχον τέλος (j = 3, 4, 5, 6)
\foreach \j [count=\r] in {3,...,6} {
  \only<\r>{
    \draw[pdpblue, very thick] (2,-0.1) -- (2,-0.3) -- (\j,-0.3) -- (\j,-0.1);
    \draw[pdpblue, thick, -latex] (\j-0.5,-1.05) -- (\j-0.5,-0.45);
    \node[text=pdpblue, below] at (\j-0.5,-1.05) {$j$};
  }
}
\end{tikzpicture}
\end{frame}

\begin{frame}[fragile]{Μέγιστο άθροισμα υποπίνακα}
    \field{Πολυπλοκότητα:}
    \begin{itemize}
        \item Τρεις φωλιασμένοι βρόχοι καθένας εκ των οποίων εκτελεί $\mathcal{O}(n)$ επαναλήψεις.
        \item Πράξεις κόστους $\mathcal{O}(1)$ εντός των βρόχων.
        \item Συνολικά $\mathcal{O}(n^3)$.
    \end{itemize}

    \begin{noteblock}{Σημείωση.}
        Ισοδύναμα, για κάθε δυνατό τέλος $j$ θα μπορούσαμε να δοκιμάζουμε δυνατή αρχή $i$, διατηρώντας την ίδια χρονική πολυπλοκότητα, κάνοντας την εξής αλλαγή. 
        \begin{cpp}
        for (int j = 1; j <= N; ++j) 
        {
            for (int i = 1; i <= j; ++i) 
            {
                // Ο υπόλοιπος κώδικας παραμένει ίδιος.
            }
        }
        \end{cpp}
    \end{noteblock}
\end{frame}

\section{Αναδρομική εξαντλητική εξέταση}
%\subsection{Διαμερισμός στοιχείων}


\pdpclosingframe               % τελευταία σελίδα με τα λογότυπα
\end{document}
